\chapter{Import and export}\label{ch:importexport}

\section{Importing}

\menu{File $\rightarrow$ Import\ldots} reads a ProTracker module (\texttt{.mod}) or a TMC song (\texttt{.tmc}, \texttt{.tm8}) and
makes a song of it. After the file is chosen the dialog of figure~\ref{fig:dialog-import} asks how to convert a module.

\dialog{dialog-import}{Importing a ProTracker module.}{11cm}

\begin{itemize}
  \item The \emph{type of RMT module}: \texttt{RMT4} (the four channels of the module on the four tracks) or \texttt{RMT8}
        (the channels spread on the two POKEYs, in the order 1, 4 / 2, 3).
  \item The treatment of the \emph{notes} (shift the octave down if the tuning of the song is too high, replace the portamento
        effects by calculated notes), of the \emph{volumes}, and the \emph{size optimizations} that are made while importing (wise loops, cutting
        the unused parts of the tracks).
\end{itemize}

When the conversion is done a dialog reports how many tracks, instruments and song lines the song has, and what the optimizations did.
Check the box and press \menu{OK}: the song is loaded, but it still needs your work. The instruments are rough: look at their envelopes and
distortions (drums, basses\ldots) and correct the tunings of the original samples.

\dialog{dialog-import-2}{The result of the import.}{9cm}

\section{Exporting}

\menu{File $\rightarrow$ Export} writes the song in a format for the Atari. The file dialog chooses the name and the format; most of the
formats then show a dialog of options. Every format is written from the data of the song as the tracker plays it.

\begin{center}
\begin{tabular}{@{}p{4.2cm}p{9.6cm}@{}}
\toprule
Format & Use \\ \midrule
RMT stripped song file (\texttt{.rmt}) & the module without the names and the unused tracks and instruments, to link in a program that uses
  the RMT player routine \\
ASM simple notation (\texttt{.asm}) & the notes as assembler data, for your own player \\
SAP-R data stream (\texttt{.sapr}) & the register writes of the POKEY, frame by frame \\
Compressed SAP-R (\texttt{.lzss}) & the same, compressed, for the LZSS player \\
SAP file + LZSS driver (\texttt{.sap}) & a file the Atari players (ASAP, the plug-ins of many players, ASMA) play directly \\
XEX Atari executable (\texttt{.xex}) & a program for the Atari that plays the song when it is started \\
Relocatable ASM for RMTPlayer (\texttt{.asm}) & the module as source with labels, to be placed in memory by your program \\
WAV audio file (\texttt{.wav}) & the sound as a 16-bit, 44.1\,kHz, stereo file, the song played once up to its loop \\
\bottomrule
\end{tabular}
\end{center}

\subsection{RMT stripped}

\dialog{dialog-export-stripped-rmt}{The export of an RMT stripped file.}{10cm}

The \emph{from address} is where the module will be in the memory of the Atari; the program shows the range it takes. The
options are \emph{SFX support} (keep the unused tracks and instruments, for sound effects played by the game),
\emph{GlobalVolumeFade} (a variable the program can use to fade the music) and \emph{No songline start} (the song always starts from
song line 0). The text box lists the \emph{RMT FEATures} the song uses: if a game uses only this song, these definitions can be pasted into
the player routine source (\texttt{rmt\_feat.a65}), which then assembles to a smaller and faster routine
(this is described in the file \texttt{optimization.txt} that comes with the program).

\subsection{ASM simple notation}

\dialog{dialog-export-simple-asm}{The export of the notes as assembler data.}{8cm}

Choose what to export (the tracks, or the whole song by song columns), how the note is written (the index \texttt{\$00}--\texttt{\$3C}, or the
frequency according to the distortion of the instrument), how the duration is written (notes only, with the special value
\texttt{XXX} in empty beats, or pairs of note and duration, in either order) and the prefix of the labels (empty for none).

\subsection{SAP, SAP-R and LZSS}

\dialog{dialog-export-sap}{The export of a SAP file.}{8cm}

A SAP file has a header with the name, the author and the date, which the dialog fills with those of the song. The field
\emph{Subsongs} lists, in hexadecimal, the song line each subsong starts from: \texttt{00} for one tune, \texttt{00 10} for two. SAP-R is the stream of
the POKEY registers; LZSS compresses it, and the SAP file contains it with the player routine of the VU-Player.

\subsection{XEX}

\dialog{dialog-export-xex}{The export of an Atari executable.}{10cm}

The program an XEX file contains shows a text on the Atari screen while it plays: up to five lines of 40 characters (the fifth shown
instead of the fourth when the \key{SHIFT} key of the Atari is held down). Below, the \emph{rasterbar} shows the CPU usage of the player
in colour, \emph{shuffle} changes its colours, and the box \emph{Automatically adjust playback speed}, with the line that says that the speed will be
adjusted to 50\,Hz on both PAL and NTSC systems, turns the adjustment of the speed on. The preview shows
the text as the Atari will draw it.

\subsection{Relocatable ASM}

\dialog{dialog-export-reloc-asm}{The export of a relocatable module.}{10cm}

The module is written as byte definitions with labels (the one of the start of the song data, and optionally separate labels for the
instruments, the tracks and the song lines, which can be placed in different places of memory). The text box shows the size of each
part and the order of the modules.

\subsection{WAV}

The WAV export has no options: it plays the song with the built-in emulation and writes the sound.

\section{Exporting from the command line}

All the formats can be exported from a script, with no window, so that the files of a release are made from the \texttt{.rmt}
source (chapter~\ref{ch:scripting}).
